\documentclass[a4paper,11pt]{article}
\usepackage[utf8]{inputenc} % allow utf-8 input
\usepackage[T1]{fontenc}    % use 8-bit T1 fonts
\usepackage{url}            % simple URL typesetting
\usepackage{booktabs}       % professional-quality tables
\usepackage{amsfonts}       % blackboard math symbols
\usepackage{nicefrac}       % compact symbols for 1/2, etc.
\usepackage{microtype}      % microtypography
\usepackage{placeins}
\usepackage{graphicx}
\usepackage[table,xcdraw]{xcolor}
\usepackage[portuguese]{babel}
\usepackage{amsmath}
\usepackage{hyphenat}
\usepackage{booktabs}
\usepackage{float}
\usepackage[sort, numbers]{natbib}

\usepackage{propp_ufjf}

% Para evitar muito espaço em braco ao longo do texto todo.
\raggedbottom
\pagenumbering{gobble} %Para remover a contagem de página, retire essa linha caso queira o número da página
% -----------------------------------------------------------------------------
\begin{document}
% -----------------------------------------------------------------------------

% RETIRAR antes de submeter
% -----------------------------------------------------------------------------
%Instruções
\includepdf[]{folha_de_rosto.pdf}
\noindent\textbf{$<$Observação: Preencha os dados solicitados sem alterar o layout desta página de rosto (capa). Na terceira página deste documento estão os itens do projeto a serem preenchidos, que deverão limitar-se a 7 folhas, excluindo a folha de rosto. {\color{red} \underline{Esta observação deve ser removida antes} \underline{de salvar o documento.}}$>$}
% -----------------------------------------------------------------------------

\vspace{1cm}

\begin{huge}
{\centering\textbf{PROPP - Pesquisa}\par}
\end{huge}

\vspace{1cm}

\begin{Large}
{\centering Dados do Projeto e do(a) Coordenador do Projeto\par}
\end{Large}

%Pagina de rosto(Capa) do projeto

\begin{table}[H]
\arrayrulecolor{gray}
\centering
\resizebox{\textwidth}{!}{%
\begin{tabular}{|
>{\columncolor[HTML]{EFEFEF}}r |p{12cm}|}
\hline
\textbf{Título do Projeto:}       &                                                                   \\ \hline
\textbf{Referência da Chamada:}  & (  ) BIC/UFJF e PIBIC/CNPq                                        \\
                                 & (  ) PIBIC/CNPq AÇÕES AFIRMATIVAS                                 \\
                                 & (  ) PIBIC-EM/CNPq                                               \\
                                 & (  ) PIBIC/FAPEMIG
                                            \\
                                 & (  ) PIBIC-JR/FAPEMIG
                                        \\
                                 & (  ) Cadastro na PROPP
                                  \\
                                 & (  ) Outra: \rule{5cm}{0.5pt}                               \\ \hline
\textbf{Coordenador do Projeto:} &                                                                   \\ \hline
\textbf{Equipe:}                 &                                                                   \\ \hline
\textbf{Endereços para contato:} & Eletrônico:                                                       \\ 
                                 & Telefônico:                                                       \\ \hline
\textbf{Unidade/Departamento:}   &                                                                   \\ \hline
\textbf{Data:}                   &                                                                   \\ \hline
\end{tabular}%
}
\end{table}

\clearpage

% -----------------------------------------------------------------------------
\section{Justificativa/Caracterização do Problema}
% -----------------------------------------------------------------------------

Faça, de forma sucinta, um relato da situação-problema abordada, citando dados ou informações significativas que possam delimitar seu contexto. Fundamente sua defesa e linha de atuação/tema (O que?; Por que?) . Apresentar revisão bibliográfica atualizada que justifique objetivamente a execução da proposta~\cite{identifier}.


% -----------------------------------------------------------------------------
\section{Objetivos}
% -----------------------------------------------------------------------------
Enuncie os objetivos a serem alcançados pelo projeto.


% -----------------------------------------------------------------------------
\section{Metodologia e Estratégias de Ação}
% -----------------------------------------------------------------------------
Explique como se pretende atingir os objetivos pretendidos. 


% -----------------------------------------------------------------------------
\section{Resultados e os impactos esperados}
% -----------------------------------------------------------------------------
Informe os resultados e/ou produtos esperados. 

% -----------------------------------------------------------------------------
\section{Cronograma}
% -----------------------------------------------------------------------------
Faça um cronograma das atividades a serem executadas no projeto. 

% -----------------------------------------------------------------------------
\section{Orçamento}
% -----------------------------------------------------------------------------

Elaborar orçamento detalhado do projeto, caso seja solicitação de recursos ou a viabilidade financeira para os outros casos.

\bibliographystyle{plain}
\bibliography{refs}

\end{document}
